\section{Packing signed intervals}
\label{sec:packing}

The sampling construction will divide the coefficient positions into
$H$ blocks, each of normalized length $1/H$. To explain the packing
required, consider one mode with coefficient $c\in\C$, a smooth cutoff
$0\le\chi\le1$ supported strictly inside its block, linear frequency
$\eta\in\R$, and nonzero quadratic curvature $\lambda\in\R$ with
$|\lambda|/H<1$. On a block centered at $x_*$ its phase, with $x=k/N$, is
\[
 N\left(\frac{\lambda}{2}(x-x_*)^2+(t+\eta)x\right).
\]
After allowing for integer frequencies in Poisson summation, a stationary
point in the cutoff support can occur only when $-t$ belongs to the circle
interval centered at
$\eta\pmod1$ of width $|\lambda|/H$. With the mode and cutoff fixed,
quadratic stationary phase gives the asymptotic bound
$|c|/\sqrt{|\lambda|}$ for its contribution divided by $\sqrt N$.
We therefore reserve total width $|\lambda|$ across the $H$ blocks for
this mode, while each block has the same peak bound. Section~\ref{sec:sampling}
will make these estimates uniform in $t$. For a real auxiliary polynomial,
opposite Fourier modes produce opposite centers and equal widths, so both
intervals must be reserved.

We now prove the geometric statement needed to separate all these
stationary intervals. Their centers must be values of specified integer
linear forms at one common point for each block. We allow relations
among three or more forms; the only independence assumption concerns
pairs of forms.

For $\eta\in\T$ and $0<\ell<1$, write
\[
 I(\eta,\ell)=\{\eta+x\pmod 1:-\ell/2\le x\le\ell/2\}
\]
for the closed interval with center $\eta$ and length $\ell$.

\begin{lemma}[Signed interval packing]
\label{lem:signed-packing}
Let $m\ge1$ and $r\ge2$ be integers, and let $a_1,\ldots,a_r\in\Z^m\setminus\{0\}$
be pairwise nonparallel.  If $w_1,\ldots,w_r>0$ satisfy
$2\sum_iw_i<1$, there exist an integer $H\ge1$ and points
$\theta_1,\ldots,\theta_H\in\T^m$ such that the $2rH$ closed intervals
\[
 I\bigl(\sigma a_i\cdot\theta_h,w_i/H\bigr),
 \qquad 1\le i\le r,\quad 1\le h\le H,\quad \sigma\in\{-1,1\},
\]
are pairwise disjoint.
\end{lemma}

We will use the following form of the Pippenger--Spencer edge-coloring
theorem \cite{PippengerSpencer1989}; the quantified statement below is
\cite[Lemma~2.1]{AlonYuster2005}.  A hypergraph here is finite and has distinct edges,
each an $r$-element subset of its vertex set.  Its degree $d(v)$ counts
edges containing $v$, and its codegree $d(u,v)$ counts edges containing
both distinct vertices $u,v$.  An edge coloring partitions the edges
into matchings, where a matching is a collection of pairwise disjoint edges.

\begin{theorem}[Pippenger--Spencer]
\label{thm:packing-edge-coloring}
For each integer $r\ge2$ and each $\gamma>0$, there exists $\beta>0$
with the following property.  If an $r$-uniform hypergraph satisfies,
for some $t>0$,
\[
 (1-\beta)t<d(v)<(1+\beta)t\quad\text{for every vertex }v,
 \qquad d(u,v)<\beta t\quad\text{for }u\ne v,
\]
then its edges can be partitioned into at most $(1+\gamma)t$ matchings.
\end{theorem}

In particular, for fixed $r$, degrees $(1+o(1))D$ and maximum
codegree $o(D)$ give an edge coloring with $(1+o(1))D$ colors.
We construct such a hypergraph whose vertices are small disjoint slots
in a half-circle.  An edge records one compatible center of each type.
A large matching will provide many compatible choices without reusing
any slot.

\begin{proof}[Proof of Lemma~\ref{lem:signed-packing}]
All constants implicit in asymptotic estimates may depend on the fixed
forms and the fixed slot lengths chosen next.  The sole parameter tending
to infinity in this proof will be a prime $q$.

\paragraph{Slot lengths and finite-field ranks.}
Choose positive proportions $\alpha_i>2w_i$ with $\sum_i\alpha_i=1$.
Approximating $T\alpha_i$ by positive odd integers as $T\to\infty$
gives fixed odd integers $L_i$ such that, with $L=\sum_iL_i$,
\begin{equation}
 w_i<\frac{L_i}{2L}\qquad(1\le i\le r).
 \label{eq:packing-width-slack}
\end{equation}
Let $b$ be the rank over $\mathbb Q$ of the matrix whose rows are the $a_i$.
Then $b\ge2$.  For every sufficiently large prime $q$, the full matrix
has rank $b$ over $\mathbb F_q$, and each pair of rows has rank two.
Indeed, it suffices to exclude the prime divisors of a nonzero
$b$-by-$b$ minor and of one nonzero two-by-two minor for each pair.
Fix such a prime for the following construction.

\paragraph{Random slots and oriented centers.}
Put $n=(q-1)/2$, $R=\lfloor\sqrt q\rfloor$, and
$U_q=\{1,\ldots,n\}$.  For $s\in\mathbb F_q$, let $|s|_q$
be the absolute value of its representative in $[-n,n]$.
Partition $U_q$ into consecutive chunks of $R$ sites, with a possible
final shorter chunk.  There are $K=\lceil n/R\rceil$ chunks.

Consider a periodic partition of the integers into consecutive slots:
each period has one slot of type $i$, consisting of $L_i$ consecutive
sites, for each $i=1,\ldots,r$.  Independently in each chunk, translate
this pattern by a uniform residue modulo $L$, and keep only slots
entirely contained in that chunk.  A kept slot is a vertex of its type.
Its middle site $c$ is an integer because $L_i$ is odd.  Give it an
oriented center $s=c$ or $s=-c$ in $\mathbb F_q$, with equal probabilities,
independently for all kept slots.  The random data in different chunks
are independent.

A chunk of length $u$ contains $u/L+O(1)$ kept slots of each type,
uniformly in its shift.  Thus in every realization the number of
vertices of each type is
\begin{equation}
 \frac nL+O(K)=\frac q{2L}+O(q/R+1).
 \label{eq:packing-vertex-count}
\end{equation}
This estimate includes the last short chunk.  Call a site interior if
its distance from both endpoints of its chunk is at least $L$.
For any nonzero $u\in\mathbb F_q$ whose folded site $|u|_q$ is interior,
the probability that $u$ is the oriented center of a vertex of a specified
type $j$ is exactly
\begin{equation}
 p=\frac1{2L}.
 \label{eq:packing-center-probability}
\end{equation}
There is exactly one shift modulo $L$ placing a type-$j$ middle site
there; the resulting slot is fully contained in the chunk, and its
orientation has the required sign with probability $1/2$.

\paragraph{Compatible tuples form a simple hypergraph.}
An edge is a tuple of one vertex of each type, with oriented centers
$s_1,\ldots,s_r$, for which the equations
\[
 a_i\cdot y=s_i\qquad(1\le i\le r)
\]
have a solution $y\in\mathbb F_q^m$.  We keep one edge per tuple,
regardless of the number of solutions.  Every compatible tuple has
exactly $q^{m-b}$ solutions, the size of the kernel of the rank-$b$
evaluation map.  This constructs a simple $r$-uniform hypergraph.

Write $C_j(u)$ for the indicator that $u$ is the oriented center of a
type-$j$ vertex.  For nonzero $s$, define the random variable
\[
 Z_{i,s}=\frac1{q^{m-1}}
 \sum_{\substack{y\in\mathbb F_q^m\\a_i\cdot y=s}}
       \prod_{j\ne i}C_j(a_j\cdot y).
\]
When the vertex of type $i$ with center $s$ exists, its degree is exactly
\begin{equation}
 d(i,s)=q^{b-1}Z_{i,s}.
 \label{eq:packing-degree}
\end{equation}
This follows by dividing the number $q^{m-1}Z_{i,s}$ of generating
vectors by $q^{m-b}$.  Likewise, two fixed evaluations of different
types leave $q^{m-2}$ possible vectors, so their codegree is at most
\begin{equation}
 q^{b-2}.
 \label{eq:packing-codegree}
\end{equation}
Two vertices of the same type have codegree zero.

We next show that nearly all vertices have nearly the same degree.
Although evaluations of several forms may be dependent, their folded
sites usually belong to different chunks.  The independent chunk
choices then supply the first and second moments we need.

\paragraph{Exceptional centers and the moment estimates.}
For a fixed type $i$, let $E_i$ contain those nonzero $s\in\mathbb F_q$
for which there exist $j<l$, both different from $i$, and a relation
\begin{equation}
 a_j+a_l=c a_i\quad\text{or}\quad a_j-a_l=c a_i,
 \qquad |cs|_q\le R.
 \label{eq:packing-exceptional-relation}
\end{equation}
Any such relation has a unique nonzero $c$: if $c=0$, the two forms
$a_j,a_l$ would be parallel over $\mathbb F_q$.  For each relation,
multiplication by $c$ is a bijection, so at most $2R+1$ residues are
excluded.  There are at most $2\binom{r-1}{2}$ relations, giving
\begin{equation}
 |E_i|=O(R).
 \label{eq:packing-exception-count}
\end{equation}

Fix a nonzero $s\notin E_i$, and sample $y$ uniformly on the affine fiber
$a_i\cdot y=s$.  Each other evaluation $a_j\cdot y$ is uniform on
$\mathbb F_q$, by pairwise rank two.  Let $Q_s$ be the chunk containing
$|s|_q$.  At most $2LK$ positive sites are noninterior; this bound also
covers a short final chunk with no interior sites.  Including their
negatives, zero, and the residues folding into $Q_s$, the probability
that a specified evaluation fails to fold into an interior site outside
$Q_s$ is at most
\[
 \tau_q=\frac{1+4LK+2R}{q}=O(q^{-1/2}).
\]

If two residues fold into one chunk, the difference of their folded
sites has absolute value at most $R$.  Consequently at least one of
their sum and difference has folded absolute value at most $R$.
For $a_j\cdot y$ and $a_l\cdot y$, each of these signed combinations
is either uniform on $\mathbb F_q$ under $a_i\cdot y=s$, or is the
constant $cs$ from a relation in
\eqref{eq:packing-exceptional-relation}.  The latter case cannot cause
a collision because $s\notin E_i$; the former has probability at most
$(2R+1)/q$.  Thus the $r-1$ evaluations fold into distinct interior
chunks different from $Q_s$, except on an $O(q^{-1/2})$ fraction of
the fiber.

For the second moment, take two independent uniform vectors $y,y'$
from the same fiber.  The preceding bounds apply within each copy.
Across copies, $a_j\cdot y$ and $a_l\cdot y'$ are independent uniform
residues for any $j,l\ne i$, including $j=l$; their sum and difference
are therefore uniform.  A union bound shows that the combined list
of $2(r-1)$ evaluations also occupies distinct interior chunks
outside $Q_s$, apart from a fraction $O(q^{-1/2})$.
More explicitly, for either $u=1$ or $u=2$ independent fiber samples,
the failure fraction is at most
\begin{equation}
 u(r-1)\tau_q+
 2\binom{u(r-1)}2\frac{2R+1}{q}=O(q^{-1/2}).
 \label{eq:packing-good-tuples}
\end{equation}
This bound uses no joint independence among the forms in one copy.

Let $V_{i,s}$ be the event $C_i(s)=1$, and condition on it whenever
it has positive probability.  The event depends only on random data
in $Q_s$.  Each tuple counted as good in
\eqref{eq:packing-good-tuples} uses distinct other chunks, where the
required center hits are independent and each has probability $p$.
The good-tuple condition itself depends only on the fiber vectors,
not on the slot choices.  Averaging over those vectors, and bounding
the contribution of bad tuples between zero and one, gives
\begin{align}
 \mathbb E[Z_{i,s}\mid V_{i,s}]&=p^{r-1}+O(q^{-1/2}),
 \label{eq:packing-first-moment}\\
 \mathbb E[Z_{i,s}^2\mid V_{i,s}]&=p^{2(r-1)}+O(q^{-1/2}).
 \label{eq:packing-second-moment}
\end{align}
The expectations are over the random slots and orientations.
The second identity follows by expanding the square as the average
over the independent pair $y,y'$.  Both error bounds are uniform in
nonexceptional $s$.  The conditioned vertex itself need not be
interior: avoiding its entire chunk makes that unnecessary.

\paragraph{Deleting the few vertices of atypical degree.}
Set $D_q=p^{r-1}q^{b-1}$ and $\varepsilon_q=q^{-1/8}$.
Equations~\eqref{eq:packing-first-moment}--\eqref{eq:packing-second-moment}
imply
\[
 \mathbb E\bigl[(Z_{i,s}-p^{r-1})^2\mid V_{i,s}\bigr]
 =O(q^{-1/2}).
\]
By Chebyshev's inequality and \eqref{eq:packing-degree}, a
nonexceptional present vertex has degree outside
$[(1-\varepsilon_q)D_q,(1+\varepsilon_q)D_q]$ with conditional
probability $O(q^{-1/4})$.  Summing over the at most $r(q-1)$ candidate
centers gives an expected $O(q^{3/4})$ such vertices.  The vertices with
centers in some $E_i$ number $O(q^{1/2})$ in every realization.
Markov's inequality therefore provides a realization with at most
$q^{7/8}$ vertices in these two classes, for all sufficiently large $q$.

Delete these vertices and all incident edges.  By
\eqref{eq:packing-codegree}, any remaining vertex loses at most
\[
 q^{7/8}q^{b-2}=O(q^{-1/8}D_q)
\]
edges.  Hence every remaining degree is $(1+O(q^{-1/8}))D_q$,
and the maximum codegree remains at most $q^{b-2}=o(D_q)$.
By \eqref{eq:packing-vertex-count}, each type still has
$(1+o(1))q/(2L)$ vertices.

\paragraph{A matching supplies the centers.}
We have obtained a hypergraph to which
Theorem~\ref{thm:packing-edge-coloring} applies.  If it has $v_q$
vertices, degree summation gives
\[
 |E|=\frac1r\sum_vd(v)=(1+o(1))\frac{v_qD_q}{r}.
\]
An edge coloring into $(1+o(1))D_q$ matchings therefore has a color
class with at least $(1-o(1))v_q/r=(1-o(1))q/(2L)$ edges.
A matching has no more edges than the number of vertices of any
one type.  We have consequently found a matching of size
\begin{equation}
 H=(1+o(1))\frac q{2L}.
 \label{eq:packing-matching-size}
\end{equation}
For each of its edges choose a generating vector $y_h\in\mathbb F_q^m$.
Using integer representatives of its coordinates, set
$\theta_h=y_h/q\pmod{\Z^m}$.  Then $a_i\cdot\theta_h$ is the
oriented center of the corresponding slot divided by $q$, modulo one.

It remains to check the geometry, including endpoints.  A slot of
type $i$ with sites $t,\ldots,t+L_i-1$ corresponds to the real interval
\[
 J=\left[\frac{t-1/2}{q},\frac{t+L_i-1/2}{q}\right].
\]
Its length is $L_i/q$, and its center is the middle site divided by
$q$.  The positive slots lie in $[1/(2q),1/2]$ and have disjoint
interiors.  Their reflections modulo one lie in
$[1/2,1-1/(2q)]$ and also have disjoint interiors.  Neighboring slots
may share an endpoint, and a slot may meet its reflection at $1/2$,
but none of these intersections reaches their interiors.

By \eqref{eq:packing-width-slack} and
\eqref{eq:packing-matching-size}, a sufficiently large prime gives
$w_i/H<L_i/q$ for every $i$.  The desired closed interval at each
signed center is therefore contained strictly inside its slot or
reflected slot.  Both signs use this same pair of slots, irrespective
of the chosen orientation.  The matching uses distinct vertices, so
all these containing slot interiors are disjoint.  This proves the
required disjointness of the closed intervals, including at $1/2$
and across zero on the circle.
\end{proof}

The prime, the matching size $H$, and the points $\theta_h$ have now
served only to establish the finite packing.  In applications we
choose them once, after the forms and widths are fixed; they impose
no restriction on a later sampling length.
